\section{What is still required for quasi-polynomial recognition}
\label{sec:research}

The new theorem removes a local obstacle: a supplied sector of raw
matching dimension three has polynomially many relevant standard
rays, with a polynomial exact enumeration and coverage proof.
It does not bound how many sectors are required.  Even listing all
maximal quadrilateral choices gives \(3^t\) possible sectors, each
containing its own coordinate faces.  Polynomial work per selected
sector therefore remains compatible with exponential total work.
The fixed corpus demonstrates the utility of the new eligibility
class, not its completeness for all unknots.

The following conditional statement records exactly what the
local result can contribute.  It is the dimension-three extension
of the preceding report's complete-family accounting statement
~\cite{MinimumEnvelopes}; its point is to expose hypotheses, not to
claim they have been established.

\begin{corollary}[Conditional composition with a complete search]
\label{res:conditional}
Suppose a deterministic procedure on an \(n\)-crossing knot diagram
visits at most \(S(n)\) source-bound states and supplied sectors,
with the following properties:
\begin{enumerate}
\item all encoded state sizes, matching systems, and normal-vector
bit lengths are polynomial in \(n\);
\item each sector has matching nullity at most three, or is reduced
to one by a checked certificate preserving all relevant candidates;
\item every state construction, failed branch, restart, transition,
terminal decision, and certificate replay together costs
\(S(n)\poly(n)\);
\item the necessary essential-disc-component predicates can be
produced and independently replayed in polynomial bit time for
each emitted vector;
\item the certified source correspondence and selection theorem
guarantee that every unknot yields a relevant disc in at least one
enumerated ray, while the terminal rules justify the other outcome.
\end{enumerate}
Then the local enumeration, coverage replay, and disc tests can be
incorporated into a complete recognizer in \(S(n)\poly(n)\) bit time.
In particular, a quasi-polynomial bound on \(S(n)\) gives a
quasi-polynomial bound on this composed algorithm.
\end{corollary}
\begin{proof}
The planar output theorem gives polynomially many rays per eligible
sector.  Proposition~\ref{pl:operations}, the determinant bit bounds,
and the coverage-verification analysis give polynomial bit cost for
their exact generation and replay.  The fourth hypothesis supplies
the remaining polynomial work per ray.  The other hypotheses bound
all visited states and justify the two recognition conclusions.
Summing these costs over the entire search yields the stated bound.
\end{proof}

The fourth hypothesis is not shorthand for expanding every normal
disc.  Classical interval-orbit methods are polynomial in compressed
input parameters and can handle exponential normal
weight~\cite{AHT2006}.  The actual component predicate, its boundary
information, and its proof representation must satisfy the same
encoded-size accounting.  Likewise, efficient verification of an
essential hierarchy does not automatically give an efficient method
for discovering it.

If a strategy has polynomially many choices per level, a proved
\(O(\log n)\) depth leads to \(n^{O(\log n)}\) states.  Depth
\(O((\log n)^2)\) instead leads to \(n^{O((\log n)^2)}\).
A bound stated in the size of the latest triangulation is
insufficient if intermediate triangulations can grow without a
bound in the original \(n\).  A progress measure must also pay for
discarded suffixes and revisits, not only for successful cuts.

\subsection{Research questions with concrete success criteria}

The most useful next work should attack an unproved hypothesis or
remove an observed bottleneck.  The following questions are arranged
approximately in that order.

\paragraph{1. Constructive coverage by a small family of sectors.}
Can the search, under the actual knot-exterior hypotheses, generate
a quasi-polynomial family guaranteed to contain an essential-disc
standard ray whenever one exists?  A useful theorem must give the
selection procedure, its encoded cost, and a reason that some
witness survives.  An existential small support for a convenient
example is insufficient: layered meridians already show why support
cardinality alone is not a reliable universal parameter.  A first
manageable target is a natural restricted class of exteriors with
a certified reduction to sectors of effective cone dimension three.

\paragraph{2. Certified effective dimension above raw nullity three.}
Many raw kernel directions never meet the nonnegative orthant.
Can the active-face certificates of the recent incoming work
identify all coordinates that vanish on the feasible cone, then
produce a basis of its exact linear span without enumerating
high-dimensional candidate rays?  Once that span has dimension at
most three, the present geometry applies in it.  The missing
implementation should certify every removed direction and preserve
all feasible points.  Compare its cost with simply handling the
original sector; do not use a heuristic feasible support as proof
of a forced zero coordinate.

\paragraph{3. Incremental exact matching kernels.}
The sparse constructor eliminates the dense zero-label allocation,
but retains dense rational elimination.  Nearby sectors differ
by a few selected types.  Can one maintain a fraction-free
factorization, cycle basis, and source-bound rank certificate under
those changes?  A successful result should bound coefficient growth
and update cost, and should remain useful when a changed label
splits a previously contracted component.  The benchmark's
source-reuse improvement supplies a baseline; updates must be
compared against that stronger baseline.

\paragraph{4. Output-sensitive exact planar geometry.}
The present half-plane construction is deliberately elementary
and has a cubic worst-case arithmetic bound.  Can power-diagram
construction and an exact segment-overlay method reduce the
geometric stage toward \(O(k^2\log k)\), plus the unavoidable
cost of writing \(tR\) native coordinates?  A proof must count
clipping by the common source polygon without charging its
boundary once per irrelevant equality.  It must retain exact
behaviour for duplicate forms, lower-dimensional minima,
coincident edges, and a degenerate feasible section.

\paragraph{5. Genuine quadratic output or a stronger geometric bound.}
The abstract grid of Section~\ref{cf:grid} has
\(\Theta(p^2)\) rays, but its potential coefficients are not a
triangulation-realization theorem.  Does there exist a growing
family of valid compact knot exteriors or solid tori with
matching nullity three and \(\Omega(k^2)\) nonlink standard
rays in one sector?  Alternatively, does the sparse face-incidence
structure force a smaller bound?  Either outcome would sharpen
the parameter that the search should optimize.  The two-cap
family settles the single-group linear constant, not this
multi-group realization question.

\paragraph{6. Structural control of the active groups.}
The explicit bound separates \(A=\sum_v(m_v-1)\) from the
cross-group term.  Is it possible to simplify a triangulation,
choose a sector, or normalize a cut so that few vertex groups
have nontrivial minimum diagrams?  One nontrivial group removes
the quadratic crossing term entirely.  The quantity to preserve
is complete witness coverage, not merely the observed number of
cells in easy examples.  A candidate invariant should include
the cost and topological legality of the simplification.

\paragraph{7. A controlled extension to matching dimension four.}
For \(d=4\), the projective section is three-dimensional and
minimum cells are convex polyhedra.  Which diagram parameters
give a useful output bound for their common refinement?  A
general all-equalities arrangement is again too coarse.
Any extension should distinguish fixed dimension from
dimension growing with \(n\), handle non-simple intersections,
and provide a checkable coverage representation with a
bit-complexity proof.  Polynomiality for each fixed dimension
would still leave the global sector-selection problem open.

\paragraph{8. Smaller certificates and a formally checked geometric core.}
The current verifier reconstructs an uncontracted source model
for independence and can therefore be slower than the producer.
Can sparse elimination be certified by local row identities,
rank witnesses, and forest-potential equations, while retaining
independence from the producing algorithm?  Another useful
project is to formalize the canonical-lift correspondence,
planar Euler bound, and area-coverage theorem in Lean or another
proof assistant.  The acceptance criterion must rule out
artificial cell subdivisions as well as missing cells; area
alone without the distinct-label and dominance conditions
does not suffice.

\paragraph{9. Source-preserving compressed topology transitions.}
After a useful surface is chosen, can cutting and continuation
retain all ordered boundary attachments, coorientations, and
pattern data with polynomial encoded size?  The hierarchy
manuscript gives relevant compressed-cutting and certification
theorems~\cite{Lackenby2026}.  The implementation milestone
should be one complete transition with independently replayed
source and target objects.  Matching an aggregate boundary
profile is insufficient when different attachments produce
different topology.

\paragraph{10. A global amortized measure compatible with restarts.}
Can local batched descent, bounded trace search, and selected
normal cuts be charged to a single well-founded quantity with
only quasi-polynomially many values?  The value must account
for regluings, unsuccessful local searches, and invalidated
later hierarchy stages.  A lexicographic genus decrease proves
termination only after every coordinate bound and every
possible length have been quantified.  This is a principal
place where a genuine global asymptotic theorem could emerge.

\paragraph{11. Integration under a measured dispatcher.}
Which exact predicate should select the old specialized
rank-two envelope, the new planar method, and the general
fallback?  The new dimension-three path has a clear coverage
advantage; the lower-dimensional controls show why a universal
replacement is unwarranted.  A useful dispatcher experiment
should retain cold-start cost, preparation reuse, coverage
replay, and the full recognizer's time to a justified verdict.
It should include both unknots and nontrivial controls and
report every capped case.

\paragraph{12. Search-independent benchmark design.}
Can a benchmark suite be assembled before tuning the selection
policy, with independently regenerated complete ray oracles
where feasible and structural families beyond them?  Hold out
source families, not just individual supports of the same
triangulation.  Report distinct surfaces as well as repeated
sector occurrences.  Measure bit lengths, number of active
cells, overlay intersections, component-query work, memory,
and complete-exhaustion time.  The present package supplies
such measurements for a local stage, but does not estimate the
frequency of these sectors in an unbiased distribution of
hard knot diagrams.

\subsection{A concrete next milestone}

The next integration step can be made precise: combine a
source-bound exterior, certified effective-face reduction,
the appropriate complete envelope enumerator, and one
compressed topology transition, with all intermediate
certificates replayed by separate checkers.  Measure the
whole step and preserve the stronger incumbent branches.
In parallel, seek a theorem bounding how often that step
must be revisited.  The first task makes the local theorem
operational; the second determines whether it advances
the global quasi-polynomial programme beyond a faster
exponential search.
